\section{Immutable Parent Pointers}

The integrity of a recursive spawning architecture depends critically on the
unforgeable linkage between each agent and its origin. Once a parent-child
relationship is established at spawn time, the resulting edge in the
derivation graph becomes a permanent, immutable record. This immutability
is not a mere operational convention; it is a first-class structural
constraint enforced at the Fact Layer. Together with the spawn authorization
mechanism at the Purpose Layer, it guarantees that every agent's chain of
custody is traceable, auditable, and tamper-evident back to the root human
at the base of the graph.

\subsection{Formal Definition of the Parent Pointer}

At spawn time, the orchestrator assigns a \texttt{ParentPointer} to the
newly created agent. This pointer is a directed, immutable reference from
the child agent $c$ to its parent agent $p$:

\begin{equation}
  \texttt{ParentPointer}(p, c) : \text{AgentID} \times \text{AgentID}
  \rightarrow \text{Pointer}
\end{equation}

The pointer is stored as an immutable record in the Fact Layer's append-only
event log. Once written, the record admits no subsequent modification or
deletion. The pointer value itself is a content-addressed identifier derived
from the parent's canonical agent record, ensuring that the reference cannot
be silently redirected to a different agent without producing a different
hash and thereby breaking the chain.

The immutability property is expressed formally as:

\begin{equation}
  \forall (p, c) \in \texttt{ParentPointer}: \;
  \neg \exists \; \texttt{modify}(\texttt{ParentPointer}(p, c), p', c')
  \quad \text{where} \quad (p', c') \neq (p, c)
\end{equation}

In other words, no operation exists that can change the target of an
established parent pointer. The only way to create a new pointer is through
a freshly authorized spawn event.

\subsection{Chain Construction}

An agent's full derivation chain is constructed by recursively applying the
parent pointer relation. We define $\texttt{Chain}(a)$ for any agent $a$ as
the set containing $a$ itself, its immediate parent, and all ancestors
recursively up to and including the root:

\begin{equation}
  \texttt{Chain}(a) = 
  \begin{cases}
    \{a\} \cup \texttt{Chain}(\text{parent}(a)) & \text{if } \text{parent}(a) \neq \texttt{null} \\
    \{a\} & \text{if } \text{parent}(a) = \texttt{null}
  \end{cases}
\end{equation}

Equivalently, this can be expressed as a union of singleton parent-pointer
edges collected along the recursive walk:

\begin{equation}
  \texttt{Chain}(a) = \{a\} \cup \bigcup_{i=1}^{n}
  \texttt{ParentPointer}(\text{parent}^i(a), \text{parent}^{i-1}(a))
\end{equation}

where $\text{parent}^i$ denotes the $i$-fold composition of the parent
function, and $n$ is the depth of $a$ in the derivation tree (the smallest
$i$ for which $\text{parent}^i(a) = \texttt{null}$).

The root human --- the originating principal at the base of the tree ---
is defined by the boundary condition:

\begin{equation}
  \text{parent}(\texttt{root\_human}) = \texttt{null}
\end{equation}

This \texttt{null} parent sentinel terminates the recursion and serves as the
only base case in the derivation graph. Every agent in the system, regardless
of how many spawn generations separate it from the root, ultimately traces
back to this single anchor point.

The depth of an agent is defined as the cardinality of its chain minus one:

\begin{equation}
  \texttt{depth}(a) = |\texttt{Chain}(a)| - 1
\end{equation}

Agents spawned directly by the root human have $\texttt{depth} = 1$, while
an agent spawned by an agent of depth $d$ has depth $d + 1$.

\subsection{Immutability Enforcement at the Fact Layer}

The Fact Layer is the foundational integrity subsystem of the BX3 framework.
It operates as an append-only ledger of cryptographically signed events.
Every spawn event, and therefore every parent-pointer assignment, is recorded
as a signed fact in this ledger. Immutability is not achieved through access
control lists or permissions policies; it is guaranteed by the ledger's
append-only nature and by the cryptographic chaining of event records.

Each fact record in the Fact Layer contains:

\begin{itemize}
  \item \texttt{event\_type}: the kind of event (e.g., \texttt{AGENT\_SPAWN})
  \item \texttt{child\_id}: the content-addressed identifier of the child agent
  \item \texttt{parent\_id}: the content-addressed identifier of the parent agent
  \item \texttt{timestamp}: a verifiable timestamp from a trusted time source
  \item \texttt{purpose\_ref}: a reference to the authorizing purpose record
        in the Purpose Layer
  \item \texttt{signature}: a cryptographic signature binding all fields above
        to the signing authority
\end{itemize}

Because each record is cryptographically self-certifying and chained to the
previous record via a hash chain, any attempt to alter a stored
\texttt{ParentPointer} would break the hash chain at that point, invalidating
the record and all subsequent records. TheFact Layer detects such tampering
through continuous hash-chain verification, which runs as a background
integrity process. Any broken chain triggers a nonconformity alert and
freezes further processing of the affected subgraph until the anomaly is
resolved through authorized governance review.

The Fact Layer also enforces immutability by refusing to accept any
modify or delete operation on a \texttt{ParentPointer} record through its
canonical API. The only sanctioned operations are:

\begin{itemize}
  \item \texttt{create\_pointer}: writes a new, immutable parent-pointer fact
  \item \texttt{read\_chain}: traverses the pointer chain for inspection or
        audit
  \item \texttt{verify\_chain}: validates the cryptographic integrity of a
        chain segment
\end{itemize}

No \texttt{update\_pointer} or \texttt{delete\_pointer} operation is
defined in the Fact Layer's authorized interface.

\subsection{Spawn Authorization at the Purpose Layer}

The Purpose Layer governs the authorization of spawn events. A spawn is not
merely a technical act of creating an agent process; it is a purposeful act
that must be justified, scoped, and authorized before the Fact Layer will
record the resulting parent pointer. This two-layer design separates the
question of whether a spawn is permitted (Purpose Layer) from the question
of how the resulting relationship is recorded (Fact Layer).

A spawn authorization is itself a signed fact stored in the Purpose Layer.
It specifies:

\begin{itemize}
  \item \texttt{spawner}: the agent or human requesting the spawn
  \item \texttt{proposed\_child\_type}: the category or class of agent to be created
  \item \texttt{authorized\_purpose}: a reference to the purpose or mission
        that justifies the spawn
  \item \texttt{scope\_constraints}: any constraints on the child agent's
        capabilities, tool access, or operational boundaries
  \item \texttt{authorization\_signature}: the signature of the authorizing
        principal
\end{itemize}

The Purpose Layer evaluates whether the proposed spawn is consistent with
the declared purpose and within the scope of the parent's authorization
grant. Only upon successful authorization does it emit an authorization
fact that the Fact Layer will accept as the basis for creating the
\texttt{ParentPointer}. This coupling ensures that parent pointers are
never created without a corresponding, audited purpose justification.

The authorization model supports delegation: a parent agent may be authorized
to spawn child agents on behalf of a higher-level purpose, subject to
scope constraints recorded in its own authorization grant. This creates
a delegation chain parallel to the parent-pointer chain, where each spawn
step carries its own purpose-level justification.

\subsection{Chain Traversal Algorithm}

Chain traversal is the core algorithmic primitive for querying the derivation
history of any agent. The following algorithm computes $\texttt{Chain}(a)$
for a given agent $a$, returning the ordered list of agents from $a$ to the
root human:

\begin{algorithm}
  \caption{Chain Traversal}
  \begin{algorithmic}[1]
    \Function{ChainTraversal}{$a$}
      \State $\text{chain} \gets []$
      \State $\text{current} \gets a$
      \Loop
        \State $\text{APPEND}(\text{chain}, \text{current})$
        \State $\text{parent\_id} \gets \texttt{FactLayer.ReadParentPointer}(\text{current})$
        \If{$\text{parent\_id} = \texttt{null}$}
          \State \Return $\text{chain}$
        \EndIf
        \State $\texttt{FactLayer.VerifyChainIntegrity}(\text{current}, \text{parent\_id})$
        \State $\text{current} \gets \text{parent\_id}$
      \EndLoop
    \EndFunction
  \end{algorithmic}
\end{algorithm}

The algorithm terminates because each iteration moves to the parent of the
current node, and the parent chain is acyclic by construction (a new
pointer is created for each spawn, but no existing pointer is ever modified
or redirected). Since the parent relation defines a strict partial order and
the root human has \texttt{null} as its parent, the loop must terminate
after at most $n$ iterations, where $n$ is the depth of the starting agent.

Two invariant properties hold for every call to \texttt{ChainTraversal}:

\begin{enumerate}
  \item \textbf{Path Completeness:} After $k$ iterations, \text{chain}
        contains exactly the first $k$ agents on the path from $a$ toward
        the root, in order.
  \item \textbf{Integrity Verification:} Each \texttt{ReadParentPointer} call
        is followed by a \texttt{VerifyChainIntegrity} call, which confirms
        that the retrieved parent pointer has not been tampered with since
        it was written. If verification fails, the algorithm aborts and
        reports a chain integrity error.
\end{enumerate}

The time complexity of \texttt{ChainTraversal} is $O(d)$, where $d$ is the
depth of the starting agent in the derivation tree. The space complexity is
also $O(d)$ for storing the resulting chain. In practice, $d$ is expected
to be bounded by a small constant (on the order of 5--10 generations) in
most operational deployments, making traversal a constant-time operation for
all practical purposes.

The ordered output of \texttt{ChainTraversal} supports several downstream
operations: audit logging, compliance verification, provenance querying,
and accountability tracing. Any consumer of agent identity or lineage
information can invoke this algorithm to reconstruct the full custodial
history of any agent in the system.

\subsection{Anchoring to the Root Human}

The root human occupies a special structural position in the derivation
graph. As the only agent whose parent is \texttt{null}, the root human
serves as the ultimate trust anchor for all descendant agents. The
immutability of the parent pointer chain ensures that no agent can,
through any sequence of spawn operations, detach itself from this anchor
or insert a false ancestor into its chain.

This design mirrors the concept of a chain of trust in public-key
infrastructure, where each certificate is valid because it is signed by
a higher-level certificate, ultimately resting on a root certificate
authority that is trusted by definition. In the BX3 recursive spawning
model, the root human is the root of trust. Every spawn event inherits its
legitimacy from the authorization chain that connects it, through a finite
sequence of parent pointers, to the root human's original activation of the
system.

The consequence is a fully auditable, cryptographically verifiable, and
tamper-evident lineage graph in which:

\begin{itemize}
  \item Every edge (parent pointer) is immutable once written.
  \item Every spawn event is authorized at the Purpose Layer.
  \item Every agent's identity is traceable to a specific, named human at
        the root.
  \item The chain admits no shortcuts, no retroactive insertion of
        ancestors, and no orphaned agents.
\end{itemize}

This forms the backbone of agentic accountability within the BX3 framework,
enabling downstream functions such as audit trails, compliance reporting,
delegation verification, and root-cause analysis to operate on a reliable,
tamper-evident data structure.
